\chapter{Phase 2: SLAM con mapa desconocido}

La escena \path{Sim_T2_Phase2_SLAM_Unknown.ttt} parte de la celda validada y
arranca con parte del entorno oculto. R1 ve cajas, paredes, una tarima móvil y
un bloque temporal a través de los sensores de proximidad dibujados como rayos.
El controlador actualiza el mapa que va armando y lo usa para activar puntos
intermedios locales.

Solo \texttt{P2\_Dynamic\_Pallet} y \texttt{P2\_Temporary\_Blocker} se mueven
por script, con trayectorias
periódicas deterministas. Las crates y los marcadores de zona desconocida quedan
fijos.

La métrica de actividad de mapa en Phase 2 cuenta rasgos detectados, celdas
ocupadas y actualizaciones SLAM. Una mayor actividad no implica necesariamente
mejor geometría; por eso la cobertura contra la referencia geométrica se revisa
aparte con exhaustividad de obstáculos y precisión de rasgos frente a esos
obstáculos.

\begin{figure}[H]
  \centering
  \begin{subfigure}[b]{0.49\textwidth}
    \centering
    \includegraphics[width=\textwidth]{coppeliasim_phase2_initial_unknown_oriented.png}
    \caption{Inicio con zona desconocida.}
  \end{subfigure}
  \hfill
  \begin{subfigure}[b]{0.49\textwidth}
    \centering
    \includegraphics[width=\textwidth]{coppeliasim_phase2_replan_dynamic_oriented.png}
    \caption{Replanificación ante obstáculo dinámico.}
  \end{subfigure}
  \caption{Capturas 3D de la escena Phase 2 ejecutada en CoppeliaSim.}
  \label{fig:phase2-coppeliasim-captures}
\end{figure}

\begin{figure}[H]
  \centering
  \begin{subfigure}[b]{0.49\textwidth}
    \centering
    \includegraphics[width=\textwidth]{A2_T2_Phase2_SS1.png}
    \caption{Zona desconocida, Bill y primer bloque no modelado.}
  \end{subfigure}
  \hfill
  \begin{subfigure}[b]{0.49\textwidth}
    \centering
    \includegraphics[width=\textwidth]{A2_T2_Phase2_SS2.png}
    \caption{Plano general de la ruta y los elementos de la celda.}
  \end{subfigure}

  \vspace{0.35em}
  \begin{subfigure}[b]{0.70\textwidth}
    \centering
    \includegraphics[width=\textwidth]{A2_T2_Phase2_SS3.png}
    \caption{Lecturas de proximidad frente al sofá.}
  \end{subfigure}
  \caption{Capturas de ejecución de Phase 2 con mapa desconocido y obstáculos en la celda.}
  \label{fig:phase2-coppeliasim-screenshots-extra}
\end{figure}

\section{Recuperación ante bloqueo local}

En Phase 2, las variantes de grid podían quedar atrapadas en estados de
recogida, con alternancia prolongada entre \texttt{AVOIDING} y
\texttt{SLAM\_WAYPOINT}. El planner trataba como bloqueo cualquier celda ocupada
cerca del segmento al objetivo, incluso cuando la celda estaba pegada al robot,
a la estantería o al punto de recogida.

El planner aplica cinco reglas para evitar ese ciclo:

\begin{itemize}
  \item Las celdas de grid necesitan varias detecciones antes de influir en
  la planificación.
  \item Los obstáculos muy cercanos al robot o al objetivo no fuerzan un punto intermedio.
  \item Los candidatos de desvío reciben penalización si salen del área útil del
  warehouse.
  \item Tras alcanzar un punto intermedio, R1 mantiene una ventana breve de ruta directa
  para evitar replanificar contra el mismo obstáculo.
  \item Si una variante de grid se queda varios segundos sin reducir distancia
  al objetivo, se activa \texttt{RECOVERY\_DIRECT}: la planificación ignora el
  mapa durante una ventana corta y conserva la evitación reactiva.
\end{itemize}

Con estas reglas los cuatro modos completan la misión en Phase 2. En la ejecución
final, Cartographer activó una recuperación del planificador; GMapping, Hector y
Kalman terminaron sin usarla.

\section{Comparación de métodos}

En Phase 2 los cuatro métodos se someten al mismo problema: arrancar con mapa
incompleto, descubrir obstáculos no modelados y seguir navegando mientras una
tarima se mueve. La comparación no se ordena por cantidad de mapa, sino por la
representación que mantiene la misión sin bloqueo:

\begin{summarybox}
\textbf{Kalman-landmark} prioriza pose estable y pocos rasgos; \textbf{GMapping}
prioriza ocupación por celda; \textbf{Hector} intenta corregir pose por
coincidencia scan--grid; \textbf{Cartographer} conserva submapas y cierres de
ciclo, pero introduce más condiciones para evaluar si una celda antigua bloquea
o no la ruta.
\end{summarybox}

\begingroup
\footnotesize
\setlength{\tabcolsep}{3pt}
\begin{longtable}{P{0.25\textwidth}P{0.16\textwidth}P{0.16\textwidth}P{0.18\textwidth}P{0.16\textwidth}}
\caption{Comparación Phase 2 con mapa desconocido y obstáculos dinámicos.}\label{tab:phase2-slam-comparison}\\
\toprule
\textbf{Métrica} & \textbf{GMapping} & \textbf{Hector} & \textbf{Cartographer} & \textbf{Kalman}\\
\midrule
\endfirsthead
\toprule
\textbf{Métrica} & \textbf{GMapping} & \textbf{Hector} & \textbf{Cartographer} & \textbf{Kalman}\\
\midrule
\endhead
Completado & Sí & Sí & Sí & Sí\\
Duración [s] & 128,00 & 129,40 & 138,65 & \textbf{124,80}\\
Ruta [m] & 15,905 & \textbf{15,735} & 16,441 & 15,752\\
RMSE pos. [m] & 0,04220 & 0,04587 & 0,03930 & \textbf{0,03185}\\
Error P95 [m] & 0,05638 & 0,07590 & 0,05773 & \textbf{0,03985}\\
Actividad de mapa [\%] & 94,851 & 94,930 & \textbf{95,327} & 94,613\\
Rasgos finales & 36 & 36 & 36 & 24\\
Exhaustividad aprox. final & 1,000 & 1,000 & 1,000 & 1,000\\
Precisión aprox. final & 0,250 & \textbf{0,333} & 0,278 & 0,250\\
Replanificaciones & 19 & 24 & 25 & \textbf{17}\\
Recuperaciones del planificador & 0 & 0 & 1 & 0\\
Batería usada [\%] & 24,576 & 24,476 & 25,916 & \textbf{24,324}\\
Submapas & 0 & 0 & \textbf{13} & 0\\
Cierres de ciclo & 0 & 0 & \textbf{4} & 0\\
\bottomrule
\end{longtable}
\endgroup

\begin{figure}[H]
  \centering
  \includegraphics[width=\textwidth]{phase2_slam_unknown_algorithm_comparison.png}
  \caption{Comparación de métodos en Phase 2.}
  \label{fig:phase2-slam-comparison}
\end{figure}

\begin{figure}[H]
  \centering
  \includegraphics[width=\textwidth]{phase2_slam_unknown_reconstruction_t60.png}
  \caption{Reconstrucción del mapa al mismo instante común, \(t=\SI{60}{s}\).}
  \label{fig:phase2-reconstruction-t60}
\end{figure}

\section{Discusión técnica}

Kalman-landmark fue el modo operativo más eficiente: terminó antes, usó menos
batería y mantuvo el menor RMSE de pose. Su mapa es compacto, suficiente para la
navegación de esta celda, aunque menos denso que las representaciones de grid.

Los modos de grid generan más rasgos de mapa y aumentan rápido la actividad de
mapeo. Esa métrica cuantifica actividad, no calidad geométrica por sí sola.
Cartographer fue la mejor grid en RMSE y actividad final, aunque
necesitó más ruta, más replanificaciones y una activación de
\texttt{RECOVERY\_DIRECT}. Hector obtuvo la mayor precisión aproximada final entre
las grids y también registró el mayor error máximo de pose. GMapping completó
la misión con menor duración que las otras grids.

Para ejecutar la misión con bajo error de pose, Kalman-landmark queda como base.
Para representar ocupación, las grids generan mapas más densos, pero su
utilidad debe evaluarse junto con error de pose, ruta, replanificaciones y
recuperaciones.
